%% =====================================================================
\chapter{The Tamper-Evident Ledger}\label{ch:ledger}

The ledger is what makes MediChain's records trustworthy. This chapter explains
the idea from first principles, then the code that implements it, then how
tampering is detected and demonstrated.

\section{The idea in one picture}

A \emph{hash} (SHA-256) turns any data into a 64-character fingerprint. Change
one character of the data and the fingerprint changes completely; it is
infeasible to find different data with the same fingerprint.

A \emph{hash chain} is a list of entries where each entry stores the fingerprint
of the previous entry, and its own fingerprint covers everything including that
link:

\begin{figure}[H]
\centering
\begin{tikzpicture}[node distance=0.9cm,
  blk/.style={draw=brand,fill=brandlight,rounded corners=2pt,align=left,font=\scriptsize,minimum width=3.5cm,inner sep=4pt}]
  \node[blk] (e0) {\textbf{seq 0} (genesis)\\ action: REGISTER\\ payloadHash: 9f2c\ldots\\ previous: \textit{null}\\ \textbf{entryHash: a41e\ldots}};
  \node[blk,right=of e0] (e1) {\textbf{seq 1}\\ action: CREATE record\\ payloadHash: 77b0\ldots\\ previous: a41e\ldots\\ \textbf{entryHash: 0c9d\ldots}};
  \node[blk,right=of e1] (e2) {\textbf{seq 2}\\ action: UPDATE record\\ payloadHash: e512\ldots\\ previous: 0c9d\ldots\\ \textbf{entryHash: 5be3\ldots}};
  \draw[arr] (e0) -- (e1);
  \draw[arr] (e1) -- (e2);
\end{tikzpicture}
\caption{Three entries of the records chain}\label{fig:chain}
\end{figure}

Two kinds of tampering become detectable:

\begin{enumerate}
  \item \textbf{Editing a ledger entry.} If someone changes entry 1's action, its
    recomputed hash no longer equals its stored \code{entryHash}. If they also
    recompute and store the new hash, entry 2's \code{previousEntryHash} no
    longer matches. They would have to rewrite every later entry.
  \item \textbf{Editing the data itself.} Each entry stores the hash of the
    protected document's content (\code{payloadHash}) at that moment. If someone
    edits the medical record directly in MongoDB, the record's current hash no
    longer equals the hash the ledger last recorded for it, and they cannot fix
    the ledger without breaking the chain (point 1).
\end{enumerate}

\section{Two chains}

\begin{itemize}
  \item \textbf{records}: every create or change of clinical and financial data
    (users registered and verified, records, file attachments, consents,
    appointments, prescriptions, bills, claims, break-glass accesses).
  \item \textbf{access}: every allow or deny decision of the zero-trust layer.
\end{itemize}

They are independent sequences in the same collection, so the high-volume
access log does not slow down reads of a record's history.

\section{The entry: \file{src\codesl{}models\codesl{}LedgerEntry\codedot{}js}}

\begin{lstlisting}[language=JavaScript,caption={LedgerEntry schema (\file{src/models/LedgerEntry.js})}]
const ledgerEntrySchema = new mongoose.Schema({
  chain: { type: String, enum: ['records', 'access'], required: true },
  seq: { type: Number, required: true }, // position in the chain, starting at 0
  entityType: { type: String, required: true },
  entityId: { type: String, default: null },
  action: { type: String, required: true },
  actorId: { type: String, default: null },
  actorRole: { type: String, default: null },
  payloadHash: { type: String, default: null }, // SHA-256 of the entity's content at this moment
  meta: { type: mongoose.Schema.Types.Mixed, default: {} },
  timestamp: { type: Date, required: true },
  previousEntryHash: { type: String, default: null }, // null only for the genesis entry
  entryHash: { type: String, required: true },
}, { versionKey: false });

// The unique (chain, seq) index makes forking impossible: two writers that read
// the same "last entry" both try to insert seq N, and the second one fails.
ledgerEntrySchema.index({ chain: 1, seq: 1 }, { unique: true });
\end{lstlisting}

The application has no route that updates or deletes ledger entries. The only
code that writes them is \code{ledgerService.append}.

\section{Writing an entry: \file{src\codesl{}services\codesl{}ledgerService\codedot{}js}}

\begin{lstlisting}[language=JavaScript,caption={computeEntryHash and appendOnce (\file{src/services/ledgerService.js})}]
function computeEntryHash(entry) {
  return sha256(canonicalJSON({
    chain: entry.chain,
    seq: entry.seq,
    entityType: entry.entityType,
    entityId: entry.entityId ?? null,
    action: entry.action,
    actorId: entry.actorId ?? null,
    actorRole: entry.actorRole ?? null,
    payloadHash: entry.payloadHash ?? null,
    meta: entry.meta || {},
    timestamp: new Date(entry.timestamp).toISOString(),
    previousEntryHash: entry.previousEntryHash ?? null,
  }));
}

async function appendOnce({ chain, entityType, entityId, action, actorId, actorRole, payload, payloadHash, meta }) {
  const last = await LedgerEntry.findOne({ chain }).sort({ seq: -1 }).select('seq entryHash').lean();
  const entry = {
    chain,
    seq: last ? last.seq + 1 : 0,
    entityType,
    entityId: entityId ? String(entityId) : null,
    action,
    actorId: actorId ? String(actorId) : null,
    actorRole: actorRole || null,
    payloadHash: payloadHash || (payload !== undefined ? hashObject(payload) : null),
    // JSON round-trip drops undefined values, which MongoDB would otherwise
    // store as null and so change the hash on read-back.
    meta: meta ? JSON.parse(JSON.stringify(meta)) : {},
    timestamp: new Date(),
    previousEntryHash: last ? last.entryHash : null,
  };
  entry.entryHash = computeEntryHash(entry);
  return LedgerEntry.create(entry);
}
\end{lstlisting}

\begin{howto}
\begin{enumerate}
  \item Read the current head of the chain (highest \code{seq}).
  \item Build the new entry with \code{seq = head.seq + 1} and
    \code{previousEntryHash = head.entryHash}.
  \item Hash the entity's content (\code{payload}) with canonical JSON
    (Section~\ref{ch:foundation}), unless a hash was supplied (files supply the
    SHA-256 of their bytes).
  \item Hash every field of the entry, store it as \code{entryHash}, insert.
\end{enumerate}
The \code{meta} round-trip is a subtle fix: an \code{undefined} value would be
hashed as ``absent'' but stored by MongoDB as \code{null}, so the recomputed hash
on read-back would differ and the chain would appear broken.
\end{howto}

\subsection{Concurrency: two requests at once}

If two requests read the same head and both insert \code{seq = N}, the chain
would fork. Two mechanisms prevent this:

\begin{lstlisting}[language=JavaScript,caption={Serialised appends with retry (\file{src/services/ledgerService.js})}]
const queues = new Map();

function enqueue(chain, task) {
  const previous = queues.get(chain) || Promise.resolve();
  const next = previous.then(task, task);
  queues.set(chain, next.catch(() => {}));
  return next;
}

function append(params) {
  return enqueue(params.chain, async () => {
    for (let attempt = 0; attempt < 5; attempt += 1) {
      try {
        return await appendOnce(params);
      } catch (error) {
        if (error.code !== 11000) throw error; // not a sequence collision
      }
    }
    throw new Error('LEDGER_APPEND_CONFLICT');
  });
}
\end{lstlisting}

\begin{itemize}
  \item \textbf{Within one process}, appends to a chain form a promise queue:
    each waits for the previous one, so they read the head one after another.
  \item \textbf{Across processes} (several API instances), the unique
    \code{(chain, seq)} index makes the second insert fail with duplicate-key
    error 11000; the writer re-reads the head and retries.
\end{itemize}

\subsection{The shortcut used by every module}

\begin{lstlisting}[language=JavaScript,caption={recordEvent (\file{src/services/ledgerService.js})}]
function recordEvent({ entityType, entity, action, actor, meta, payloadHash }) {
  return append({
    chain: 'records',
    entityType,
    entityId: entity._id,
    action,
    actorId: actor ? actor.userId || actor._id : null,
    actorRole: actor ? actor.role : 'system',
    payload: typeof entity.ledgerPayload === 'function' ? entity.ledgerPayload() : undefined,
    payloadHash,
    meta,
  });
}
\end{lstlisting}

Each protected model defines \code{ledgerPayload()}: the fields whose integrity
matters. For a medical record:

\begin{lstlisting}[language=JavaScript,caption={What the ledger protects for a record (\file{src/models/MedicalRecord.js})}]
medicalRecordSchema.methods.ledgerPayload = function ledgerPayload() {
  return {
    patientId: this.patientId.toString(),
    doctorId: this.doctorId ? this.doctorId.toString() : null,
    recordType: this.recordType,
    title: this.title,
    description: this.description,
    diagnosis: this.diagnosis,
    findings: this.findings,
    attachments: this.attachments.map((a) => ({ fileId: a.fileId, sha256: a.sha256 })),
    version: this.version,
  };
};
\end{lstlisting}

Bookkeeping fields such as \code{updatedAt} or \code{ledgerEntryId} are left out
on purpose: they change for harmless reasons and would cause false alarms. The
same pattern exists on \code{Appointment}, \code{AccessGrant},
\code{Prescription}, \code{Bill} and \code{Claim}.

\section{Verifying}

\subsection{Walking a chain}

\begin{lstlisting}[language=JavaScript,caption={verifyChain (\file{src/services/ledgerService.js})}]
async function verifyChain(chain) {
  const cursor = LedgerEntry.find({ chain }).sort({ seq: 1 }).lean().cursor();
  let previous = null;
  let checked = 0;

  for await (const entry of cursor) {
    const expectedSeq = previous ? previous.seq + 1 : 0;
    if (entry.seq !== expectedSeq) {
      return { chain, valid: false, entriesChecked: checked, brokenAtSeq: expectedSeq, problem: 'missing_entry' };
    }
    if (entry.previousEntryHash !== (previous ? previous.entryHash : null)) {
      return { chain, valid: false, entriesChecked: checked, brokenAtSeq: entry.seq, problem: 'broken_link' };
    }
    if (computeEntryHash(entry) !== entry.entryHash) {
      return { chain, valid: false, entriesChecked: checked, brokenAtSeq: entry.seq, problem: 'entry_modified' };
    }
    previous = entry;
    checked += 1;
  }

  return { chain, valid: true, entriesChecked: checked, headHash: previous ? previous.entryHash : null };
}
\end{lstlisting}

A database cursor streams entries one by one, so even a large chain is verified
in constant memory. The three checks detect a deleted entry (a gap in
\code{seq}), a re-ordered or re-linked entry, and an edited entry.

\subsection{Checking a document against the ledger}

\begin{lstlisting}[language=JavaScript,caption={verifyEntity (\file{src/services/ledgerService.js})}]
async function verifyEntity(entityType, entity) {
  const latest = await LedgerEntry.findOne({
    chain: 'records', entityType, entityId: String(entity._id), payloadHash: { $ne: null },
    action: { $nin: ['FILE_ATTACHED', 'VIEW', 'BREAK_GLASS_ACCESS'] },
  }).sort({ seq: -1 }).lean();

  if (!latest) {
    return { entityId: String(entity._id), matches: false, reason: 'no_ledger_entry' };
  }
  const currentHash = hashObject(entity.ledgerPayload());
  return {
    entityId: String(entity._id),
    matches: currentHash === latest.payloadHash,
    currentHash,
    ledgerHash: latest.payloadHash,
    ledgerSeq: latest.seq,
  };
}
\end{lstlisting}

The most recent content-bearing entry for the document holds the hash of what
the document \emph{should} contain. Recompute the document's hash now; if they
differ, it was changed without going through the API. A document with no ledger
entry at all (inserted directly into the database) is also reported.

\subsection{The full audit: \code{GET \codesl{}api\codesl{}ledger\codesl{}verify}}

\begin{lstlisting}[language=JavaScript,caption={Full integrity audit (\file{src/modules/ledger/ledger.routes.js})}]
const PROTECTED = [
  { entityType: 'record', model: MedicalRecord, label: 'Medical records', describe: (d) => d.title },
  { entityType: 'prescription', model: Prescription, label: 'Prescriptions', describe: (d) => `${d.medication} ${d.dosage}` },
  { entityType: 'consent', model: AccessGrant, label: 'Consents', describe: (d) => `consent ${d.status}` },
  { entityType: 'appointment', model: Appointment, label: 'Appointments', describe: (d) => d.startTime.toISOString() },
  { entityType: 'claim', model: Claim, label: 'Insurance claims', describe: (d) => d.claimNumber },
  { entityType: 'bill', model: Bill, label: 'Bills', describe: (d) => d.billNumber },
];

router.get('/verify', asyncHandler(async (req, res) => {
  const started = Date.now();
  const chains = await Promise.all(['records', 'access'].map((c) => ledgerService.verifyChain(c)));

  const entities = [];
  for (const { entityType, model, label, describe } of PROTECTED) {
    const docs = await model.find();
    const results = await Promise.all(docs.map((doc) => ledgerService.verifyEntity(entityType, doc)));
    const tampered = results.map((r, i) => ({ ...r, name: describe(docs[i]) })).filter((r) => !r.matches);
    entities.push({
      entityType,
      label,
      checked: results.length,
      tampered: tampered.map((t) => ({ entityId: t.entityId, name: t.name, reason: t.reason || 'content_changed', ledgerSeq: t.ledgerSeq })),
    });
  }

  const valid = chains.every((c) => c.valid) && entities.every((e) => e.tampered.length === 0);
  res.json({ success: true, data: { valid, chains, entities, checkedAt: new Date(), durationMs: Date.now() - started } });
}));
\end{lstlisting}

This is what the admin's \textbf{Verify ledger} button calls, and what
\code{npm run ledger:verify} runs from the command line. For one record,
\code{GET \codesl{}api\codesl{}records\codesl{}:id\codesl{}verify} additionally re-reads every attached file from
GridFS and compares the SHA-256 of its bytes with the hash in the ledger
(Chapter~\ref{ch:records}).

\section{Demonstrating tamper detection}

\file{scripts\codesl{}tamper-demo\codedot{}js} plays a malicious insider with
direct database access:

\begin{lstlisting}[language=JavaScript,caption={The insider edit (\file{scripts/tamper-demo.js})}]
  const record = await MedicalRecord.findOne({ title: 'Stage 1 essential hypertension' });
  if (!record) throw new Error('Seeded record not found. Run "npm run seed -- --reset" first.');
  fs.writeFileSync(STATE, JSON.stringify({ recordId: record._id.toString(), original: record.diagnosis }));
  // updateOne on the raw collection: no API, no validation, no ledger entry. What an insider with DB access could do.
  await MedicalRecord.collection.updateOne({ _id: record._id }, { $set: { diagnosis: 'No hypertension. Patient healthy.' } });
\end{lstlisting}

\begin{lstlisting}[style=shell,caption={Running the demonstration}]
npm run demo:tamper            # change the diagnosis directly in MongoDB
npm run ledger:verify          # -> record ... TAMPERED: <id>  ...  TAMPERING DETECTED (exit 1)
npm run demo:tamper -- --restore
npm run ledger:verify          # -> LEDGER VALID
npm run demo:tamper -- --ledger    # rewrite a ledger entry -> chain "records": BROKEN (entry_modified)
\end{lstlisting}

The same result appears in the admin dashboard (Ledger $\to$ Verify ledger) and
on the patient's record (Verify integrity). The integration test in
\file{tests\codesl{}integration\codesl{}records\codedot{}test\codedot{}js} automates
exactly this scenario:

\begin{lstlisting}[language=JavaScript,caption={Tamper detection test (\file{tests/integration/records.test.js})}]
  test('editing a record directly in MongoDB (bypassing the API) is detected', async () => {
    const w = await world();
    await grantAccess(w);
    const created = await newRecord(w.doctor, w.patientUser._id);
    const id = created.body.data._id;

    await MedicalRecord.updateOne({ _id: id }, { $set: { diagnosis: 'Healthy, no hypertension' } });

    const verify = await as(w.doctor).get(`/api/records/${id}/verify`);
    expect(verify.body.data.intact).toBe(false);
    expect(verify.body.data.content.matches).toBe(false);

    const audit = await as(w.admin).get('/api/ledger/verify');
    expect(audit.body.data.valid).toBe(false);
    expect(audit.body.data.chains.every((c) => c.valid)).toBe(true); // the chain itself is intact
    const recordCheck = audit.body.data.entities.find((e) => e.entityType === 'record');
    expect(recordCheck.tampered.map((t) => t.entityId)).toEqual([id]);
  });
\end{lstlisting}

\section{Limits and how to close them}\label{sec:ledger-limits}

A hash chain proves that history is \emph{internally consistent}. An attacker
with full write access to MongoDB \emph{and} the patience to recompute every
hash after the edited point could produce a new, consistent chain. Two
inexpensive measures defeat that:

\begin{itemize}
  \item \textbf{Anchoring.} Periodically publish the head hash of each chain
    (\code{GET \codesl{}api\codesl{}ledger\codesl{}stats} returns it) somewhere the attacker cannot
    change: a signed email to auditors, a public timestamping service, or a
    public blockchain transaction. A rewritten chain will not reproduce a
    published head hash.
  \item \textbf{Separation of duties.} Run the API with a MongoDB user that may
    insert into \code{ledgerentries} but not update or delete.
\end{itemize}

Neither is implemented yet; both are listed in Chapter~\ref{ch:demo}. Moving to
Hyperledger Fabric later would replace \code{ledgerService} only: every caller
already goes through \code{append} and \code{recordEvent}.
